\documentclass{ieeeaccess}
\usepackage{cite}
\usepackage{amsmath,amssymb,amsfonts}
\usepackage{graphicx}
\usepackage{textcomp}
\usepackage{booktabs}
\usepackage{float}
\usepackage{algorithm}
\usepackage{algpseudocode}
\usepackage{makecell}
\usepackage{url}

\usepackage{bm}
\makeatletter
\AtBeginDocument{\DeclareMathVersion{bold}
\SetSymbolFont{operators}{bold}{T1}{times}{b}{n}
\SetSymbolFont{NewLetters}{bold}{T1}{times}{b}{it}
\SetMathAlphabet{\mathrm}{bold}{T1}{times}{b}{n}
\SetMathAlphabet{\mathit}{bold}{T1}{times}{b}{it}
\SetMathAlphabet{\mathbf}{bold}{T1}{times}{b}{n}
\SetMathAlphabet{\mathtt}{bold}{OT1}{pcr}{b}{n}
\SetSymbolFont{symbols}{bold}{OMS}{cmsy}{b}{n}
\renewcommand\boldmath{\@nomath\boldmath\mathversion{bold}}}
\makeatother

\def\BibTeX{{\rm B\kern-.05em{\sc i\kern-.025em b}\kern-.08em
    T\kern-.1667em\lower.7ex\hbox{E}\kern-.125emX}}

\renewcommand{\thetable}{\Roman{table}}

%Your document starts from here ___________________________________________________
\begin{document}
\history{Date of publication xxxx 00, 0000, date of current version xxxx 00, 0000.}
\doi{10.1109/ACCESS.2024.0429000}

\title{Mitigating Class Imbalance in Pump-and-Dump Detection: A Comparative Analysis of Imbalance-Aware Algorithms}
\author{\uppercase{M. Mironov}\authorrefmark{1},
\uppercase{D. Bogutsky}\authorrefmark{2}, and P. \uppercase{Lukianchenko}\authorrefmark{3}}

\address[1]{National Research University Higher School of Economics, Faculty of Computer Science, Laboratory of AI in Mathematical Finance (e-mail: mvmironov@hse.ru)}

\address[2]{National Research University Higher School of Economics, Faculty of Computer Science, Laboratory of AI in Mathematical Finance (e-mail: dbogucki@hse.ru)}

\address[3]{National Research University Higher School of Economics, Faculty of Computer Science, Laboratory of AI in Mathematical Finance (e-mail: plukyanchenko@hse.ru)}

\tfootnote{
This work was supported by the grant for research centers in the field of AI provided by the Ministry of Economic Development of the Russian Federation in accordance with the agreement 000000C313925P4E0002 and the agreement with HSE University № 139-15-2025-009}

\markboth
{Mironov, Bogutsky, Lukianchenko \headeretal: Mitigating Class Imbalance in Pump-and-Dump Detection}
{Mironov, Bogutsky, Lukianchenko \headeretal: Mitigating Class Imbalance in Pump-and-Dump Detection}

\corresp{Corresponding author: D. Bogutsky (e-mail: dbogucki@hse.ru).}

\begin{abstract}
The cryptocurrency market has been subject to various forms of exploitation, with pump-and-dump schemes being a persistent problem that can lead to significant price volatility and financial losses. Previous studies have focused on predicting attack targets using machine learning approaches; however, these methods are often compromised by class imbalance and biases inherent in historical data. To address this limitation, we present a novel framework for predicting manipulation targets.
We employ two-step bias-minimizing data normalization techniques to mitigate the effects of class imbalance and temporal heterogeneity. We also propose several innovative machine learning approaches including imbalance-aware hyperparameter optimization and ranking algorithms.
Validation Top@k\%AUC selects the CatBoost model with Top@k\%AUC early stopping ($0.667$). Frozen before test evaluation, this model attains test Top@1 of $0.150$, Top@10 of $0.563$, and Top@k\%AUC of $0.674$ on 80 held-out events with an imbalance ratio of approximately $1{:}287$; tuned logistic regression has the highest descriptive test Top@k\%AUC ($0.697$).
We use only the validation-selected CatBoost model for the portfolio backtest and assess it under realistic, slippage-aware execution.
\end{abstract}

\begin{keywords}
Classification, Imbalanced Datasets, Learning to rank, Machine Learning, Market Manipulation, Market Microstructure, Portfolio strategy
\end{keywords}

\titlepgskip=-21pt

\maketitle

\section{Introduction}
\label{sec:introduction}
\subsection{What are Pumps and Dumps}

\PARstart{P}{ump-and-dump} schemes (P\&Ds) are a type of fraudulent activity that involves artificially inflating the price of a cryptocurrency through coordinated buying efforts, followed by a rapid selling to unsuspecting investors. This phenomenon is often characterized by an abrupt and significant increase in trading volume, followed by a subsequent decline in price as the original buyers liquidate their positions.

These schemes typically involve experienced traders and market manipulators who orchestrate the operation through various channels, including online forums and social media platforms such as Telegram channels. The manipulation is designed to create a false sense of urgency and scarcity, which attracts inexperienced investors who may be unaware of the true nature of the event.

A typical pump-and-dump scenario can be observed in Fig.~\ref{fig:pump_and_dump_example}, which shows a clear example of this phenomenon. The figure illustrates price movements associated with a specific token VIB traded against BTC on Binance exchange, highlighting the abrupt spike in trading activity following the announcement of the token's release. This sudden increase in volume is often followed by a rapid decline in price as insiders sell their positions to new buyers.

\begin{figure}[h!]
    \centering
    \includegraphics[scale=0.5]{images/pump_and_dump_example.png}
    \caption{VIBBTC pump and dump manipulation carried out on Binance}
    \label{fig:pump_and_dump_example}
\end{figure}

\subsection{Related Work}

The phenomenon of pump-and-dump schemes has been extensively studied in the context of cryptocurrencies, with research efforts spanning multiple approaches and methodologies.

One of the pioneering studies on this topic was \cite{kamps_moon_2018}, which introduced a comprehensive framework for understanding pump-and-dump operations. The authors highlighted the targeting of small market cap cryptocurrencies due to their lower trading volumes, utilizing data from 5 prominent crypto exchanges: Binance, Bittrex, Kraken, Kucoin, and Lbank. Their analysis covered pumps history of 1.5 years, with price and volume data being used as the primary sources. Subsequent research efforts have focused on leveraging machine learning (ML) algorithms to enhance P\&D detection accuracy. La Morgia et al.'s work \cite{la_morgia_pump_2020}, \cite{la2023doge} employed random forest algorithms, achieving improved results compared to traditional rule-based approaches. The authors in \cite{chadalapaka_crypto_2022} further increased the effectiveness of La Morgia's work by incorporating Deep Neural Networks (LSTMs and Transformers) as anomaly detection algorithms.
More recently, in \cite{econometrics11030022} the authors built on these findings by studying a similar pump dataset between 2021 and 2022. They developed an innovative approach that employed resampling methods, such as SMOTE, combined with price and volume features and random forest classification to detect pump events 1 hour in advance.

In addition to pump detection, other research has focused on predicting the maximum price move during known pump events. Nghiem et al. \cite{nghiem_2021} used market data and social sentiment with deep neural networks (LSTMs and CNNs) and achieved promising results. However, their analysis highlighted a potential limitation of using sentiment signals: historical bias may be introduced due to the lack of precise snapshots of historical social data.

A relevant example of work addressing the problem of predicting price manipulation targets is the study \cite{pd_anatomy}. This research proposed an approach using random forest algorithms to identify potential pump and dump schemes among listed tokens. The authors analyzed a dataset of 180 pump events from the Cryptopia exchange, incorporating features such as price and volume data, global market capitalization data from Coinmarketcap, and token ratings by exchange. To enhance model robustness, the study employed filtering techniques to exclude pumps with suspicious announcement times that deviated significantly from full hours. Moreover, the authors validated their classification model through backtesting a portfolio strategy based on predicted pump targets. However, it is worth noting that the employed trading strategy relied on unrealistic assumptions, such as guaranteed profits of half the maximum price move during the pump and zero losses for false positive predictions.

A related study was conducted in \cite{hu2023sequence}, where authors enhanced their approach by incorporating channel embeddings into their model. They also utilized standard price and volume metrics, as well as external data sources such as market capitalization and social media information (e.g., Twitter followers, Reddit subscribers, Alexa rank). The analysis was performed on a dataset comprising 445 pump announcement events. However, due to the presence of duplicate pump events across multiple channels, the actual number of unique pump events studied is smaller than expected. Furthermore, the use of social data from different sources raises concerns about historical bias, as these metrics may have changed after the fact, potentially affecting the accuracy of the results.

Regarding market manipulation in a broad sense, several papers have also been published recently, for example \cite{cong_washtrading_2023}, which studied the manipulation of trading volumes, also known as wash trading.

With the increasing popularity of ML applications in finance, in recent years we have seen many articles on predictions using the market microstructure, for example \cite{albers_frag_2022} and \cite{ntakaris_midprice_2022}. Promising results have also been published with applications of ranking models to finance \cite{poh_ltr}.  Taking into account the rich accumulated experience in related areas, we aim to complement existing work with new factors based on the market microstructure.

In the broader ML literature, several families of techniques have been proposed for handling class imbalance beyond resampling, and each maps imperfectly onto the P\&D target-prediction problem. Cost-sensitive learning~\cite{elkan2001cost} re-weights the loss so that misclassifying the minority class is more expensive; it does so without perturbing the feature space, which, as we show empirically in Section~\ref{section:results}, is a decisive advantage when features are already cross-sectionally normalized. Focal loss~\cite{lin2017focal}, originally proposed for dense object detection, down-weights easy examples and concentrates the gradient on hard, informative ones; it has since been adopted for imbalanced tabular classification, but remains most effective when combined with deep models and large minority classes, conditions that do not apply to our 290 positive training events. Synthetic oversampling methods such as SMOTE~\cite{chawla2002smote} have known pathologies in high-dimensional settings~\cite{blagus2013smote}, which we confirm in Section~\ref{section:results}.

Anomaly detection provides a complementary formulation that does not require balanced training data. Isolation Forest~\cite{liu2008isolation} and related one-class approaches have been applied to fraud and market-abuse detection when positive examples are essentially unavailable. In our setting, however, positive examples \emph{are} available (at least in limited numbers) and the task is explicitly \emph{cross-sectional}: we do not ask ``is this observation anomalous in isolation?'' but rather ``which asset in this cross-section is the most likely manipulation target?''. This ranking-within-cross-section structure makes the problem closer to learning-to-rank than to unsupervised anomaly detection, and motivates the imbalance-aware ranking metric (Top@k\%AUC) that we adopt as the training target throughout this paper. To the best of our knowledge, no prior work on P\&D target prediction has combined cross-sectional panel normalization, microstructure features, and an imbalance-aware ranking objective explicitly tailored to the 1-positive-vs.-many-negatives structure of each cross-section.


This paper builds upon previous work in \cite{pd_anatomy} and \cite{hu2023sequence}, which focused on predicting manipulation targets among traded assets. We identify several critical limitations in data preparation that limit model reliability and propose innovative solutions to address these issues. Specifically, we introduce novel microstructure-based features leveraging insights from areas of market prediction \cite{albers_frag_2022}, \cite{ntakaris_midprice_2022} and wash trading detection \cite{cong_washtrading_2023}. Moreover, we explore the use of boosted trees, as implemented in CatBoost \cite{catboost}, which has been shown to outperform other algorithms on tabular data (e.g., \cite{grinsztajn_2022}), and we compare their performance to the random forest model introduced in \cite{random_forest_2001}.

To further improve the classification models' performance, we also perform a cross-sectional normalization method for panel data. Moreover, we formulate a specialized metric tailored to optimize the classification model for highly unbalanced datasets, which was used for model training and validation.

Table~\ref{table:comparison_prior_work} summarizes the key methodological differences between our framework and the most closely related prior studies.

\begin{table*}[!t]
    \centering
    \small
    \renewcommand{\arraystretch}{1.3}
    \setlength{\tabcolsep}{5pt}
    \caption{Comparison of the proposed framework with closely related prior work on P\&D target prediction}
    \begin{tabular}{p{2.8cm}p{2.2cm}p{2cm}p{2.2cm}p{2.5cm}p{2.5cm}}
        \toprule
        \textbf{Study} & \textbf{Exchange} & \textbf{\# Events} & \textbf{Model} & \textbf{Imbalance Handling} & \textbf{Backtesting} \\
        \midrule
        Xu \& Livshits~\cite{pd_anatomy} & Cryptopia & 180 & Random forest & None & Unrealistic assumptions \\
        Hu et al.~\cite{hu2023sequence} & Multiple & 445 & Seq. model + channel embed. & None & None \\
        Fantazzini \& Xiao~\cite{econometrics11030022} & Binance & Varies & Random forest & SMOTE & None \\
        La~Morgia et al.~\cite{la_morgia_pump_2020} & Multiple & 1111 & Random forest & None & None \\
        \textbf{Ours} & Binance & 473 & CatBoost, LR, RF, Ranker & Cross-sectional norm., Top@k\%AUC opt., class weighting & Impact-aware execution \\
        \bottomrule
    \end{tabular}
    \label{table:comparison_prior_work}
\end{table*}

\section{Data Collection and Preprocessing}

\subsection{Labeled P\&Ds data}
\label{section:pump_and_dump_data}

The released event manifest, \texttt{resources/pumps.json}, contains 498 unique announced P\&D targets on Binance spot BTC pairs between January~3, 2018 and August~14, 2022. Each record supplies the target pair, exchange, and UTC announcement timestamp. These announcement labels are the fixed input to the empirical pipeline; no event is retained or discarded according to its subsequent return.

\paragraph{Channel identification.}
We identified eight candidate Telegram channels through a two-step procedure. First, we searched Telegram directly using the keywords ``crypto pump,'' ``pump signal,'' ``pump and dump,'' and their localized equivalents, retaining public channels whose stated purpose was the coordination of pump events. Second, we cross-referenced the resulting list with channels named or cited in prior literature~\cite{la_morgia_pump_2020, pd_anatomy, econometrics11030022}. For each channel, we exported the full message history (JSON format) using the Telegram Desktop client's built-in export functionality. The channel names cannot always be redistributed, but the timestamps and tickers of the retained events are in the public event file.

\paragraph{Labeling procedure.}
The operational label record contains one row per \{ticker, exchange, announcement timestamp\}; there are no exact or ticker--timestamp duplicates. The repository does not contain Telegram message text or message URLs, so \texttt{resources/pumps.json} supports exact computational reproduction of the experiments but not an independent reconstruction of the original social-media annotation decisions.

\paragraph{Inclusion criteria.}
A manifest event enters a model sample only if \emph{all} of the following data-availability conditions hold:
\begin{itemize}
    \item the target is a Binance spot pair quoted in BTC;
    \item the target has raw trade data in the exact past-only eligibility interval $[T-15\text{min}-1\text{d},T-15\text{min})$;
    \item a feature cross-section can be materialized from the local Binance archive.
\end{itemize}

\paragraph{Exclusion criteria.}
No post-announcement price or volume threshold is used for sample inclusion. If a manifest event cannot be materialized because the required raw target history is absent, \texttt{create\_dataset()} records the event in \texttt{resources/dropped\_pumps.json}. This file is a reproducible missing-feature artifact; it does not encode social-media rejection reasons.

We performed the following preprocessing steps:
\begin{enumerate}
    \item parse and validate every manifest record;
    \item discover candidate BTC pairs only from date partitions intersecting the past eligibility interval;
    \item inspect actual trade timestamps rather than date-partition membership, and require at least one trade in $[T-15\text{min}-1\text{d},T-15\text{min})$ for each candidate and the announced target;
    \item materialize features and split complete event cross-sections chronologically.
\end{enumerate}

The resulting dataset contains 473 unique market manipulation events. Restricting the analysis to Binance BTC pairs is a deliberate design choice with known trade-offs: (i)~we lose events that targeted USDT-quoted pairs, tokens listed only on smaller venues (e.g., KuCoin, Gate, LBank, Cryptopia before its 2019 shutdown), and decentralized-exchange pumps; (ii)~in exchange, we obtain full trade-level history for every pair in the cross-section, which is the minimum requirement for computing slippage, flow-imbalance and power-law features consistently across the whole sample; (iii)~the restriction is aligned with the most comparable prior work on Binance P\&D target prediction~\cite{econometrics11030022} and on microstructure-based detection~\cite{la_morgia_pump_2020}, so our results remain directly comparable. We further discuss the limits of external validity in Section~\ref{sec:limitations}. Our data presents a substantial challenge, with P\&D events occurring relatively infrequently compared to normal trading activity. During our study period, the average eligible cross-section contains 195.05 instruments quoted in BTC. The historical dynamics of the number of listings is shown in Fig.~\ref{fig:cross-section_size_over_time}.

\begin{figure}[!t]
    \centering
    \includegraphics[scale=0.5]{images/cs_size.jpeg}
    \caption{Binance has been increasing the number of tickers traded against BTC by listing more tokens. 
    %This tendency was likely to remain up until the end 2022 when more regulation came in with much stricter listing requirements.
    }
    \label{fig:cross-section_size_over_time}
\end{figure}

We split the data into training, validation and testing samples by 2020-09-01 and 2021-05-01 respectively. Table~\ref{table:train_val_test_stats} describes the resulting datasets after the split: 

\begin{table}[!t]
    \centering
    \setlength{\tabcolsep}{8pt} % Adjust the column spacing
    \renewcommand{\arraystretch}{1.2} % Adjust the row spacing
    
        \begin{tabular}{lrrr}
            \toprule
             & Train & Validation & Test \\
            \midrule
                Positive & 290 & 103 & 80 \\
                Negative & 44578 & 24274 & 22935 \\
                Total & 44868 & 24377 & 23015 \\
                Avg. Cross-section Size & 154.72 & 236.67 & 287.69 \\
            \bottomrule
        \end{tabular}
        \caption{P\&D Train-Validation-Test samples statistics}
     \label{table:train_val_test_stats}
\end{table}

\subsection{Market data}

To enable comparison with existing research, we selected Binance as the sole exchange platform. We use Binance's public historical trade archive for every BTC-quoted pair needed by the 2018--2022 event manifest and its 30-day feature lookbacks.

A representative example of our collected data is presented in Table~\ref{table:collected_tick_data}.

\begin{table*}[!t]
    \centering
    \begin{tabular}{lrrrrr}
        \toprule
         & price & qty & time & isBuyerMaker & quote \\
        \midrule
        0 & 0.000002 & 2437.00 & 2022-07-15 15:00:02.422000 & True & 0.004240 \\
        1 & 0.000002 & 275.00& 2022-07-15 15:00:02.434000 & False & 0.000479 \\
        2 & 0.000002 & 1125.00 & 2022-07-15 15:00:22.884000 & True & 0.001946 \\
        3 & 0.000002 & 941.00 & 2022-07-15 15:00:52.277000 & True & 0.001628 \\
        4 & 0.000002 & 184.00 & 2022-07-15 15:00:52.277000 & True & 0.000318 \\
        \bottomrule
    \end{tabular}
    \caption{Structure of collected data: \textit{price} shows the price that the order was filled at, \textit{qty} is the amount of base asset traded, \textit{isBuyerMaker} indicates if the buyer's order related to this transaction was already in the orderbook: if True, the trade is classified as a Sell event and if False the trade is characterized as a Buy event.}
    \label{table:collected_tick_data}
\end{table*}


\subsection{Feature engineering}

We collected trade-level data as compressed daily archives for the required Binance pairs and dates. Because trading activity varies substantially across tickers and over time, the raw event streams are irregular and cannot be used directly as fixed-size inputs to standard machine-learning models. We therefore implemented a feature-generation pipeline that (i) decompresses and loads the relevant archives, (ii) partitions data into time-aligned windows, (iii) computes microstructure features summarizing pre-announcement order flow, and (iv) aggregates ticker-level features into an event-level cross-section associated with each pump. The pipeline is:

\begin{enumerate}
    \item Download trade-level data from Binance's public daily archive.
    \item For each pump event, construct the cross-section from tickers with at least one trade in the exact interval $[T-15\text{min}-1\text{d},T-15\text{min})$. Partition membership alone is insufficient because a daily partition can contain trades later than the decision time.
    \item For each eligible ticker, load 30 days of history ending at the decision cutoff and compute features over predefined half-open lookback windows $[T-15\text{min}-x,T-15\text{min})$. Hourly normalizer bars are re-anchored to the cutoff so every included bar is strictly pre-decision.
    \item Combine features for all tickers into a single event-level dataset and assign labels by adding \textit{is\_pumped}, which equals 1 for the pumped ticker and 0 otherwise.
\end{enumerate}

This modular pipeline allowed fast iteration over alternative feature sets. Implementing it efficiently required substantial effort due to the scale of the data; the complete audited regeneration of 498 manifest events finished in approximately 28 minutes with 16 worker processes. The implementation and full instructions to reproduce our results are available in \cite{code_github}. We group features into several categories, described below.

\subsubsection{Imbalance features}

Volume imbalance measures which side is more prevalent in trade data. It is the ratio of signed notional in the quote asset (negative for seller-initiated trades and positive for buyer-initiated trades) to total absolute quote notional. Values near $1$ indicate predominantly buyer-initiated flow and values near $-1$ predominantly seller-initiated flow.

\begin{equation}
    \text{Volume Imbalance} = \frac{\sum_{t \in \mathcal{T}}s_t p_t q_t}{\sum_{t \in \mathcal{T}}p_t q_t},
    \label{eq:volume_imbalance}
\end{equation}
where $p_tq_t$ is quote notional and $s_t\in\{-1,1\}$ is the taker-side sign at time $t$.

\subsubsection{Quote slippages}

We also created features based on slippage. Since we do not have information about order types, we cannot determine whether a trade was executed via a market or limit order; therefore, we treat slippage as a useful proxy for measuring market participants’ impatience. We follow the methodology used in prior work by grouping ticks by execution time. Given that we deal with highly illiquid markets—where the time between trades can be as long as 15 minutes—we can reasonably assume that aggregated ticks correspond to a single order. This allows us to estimate the price impact of the aggregated order. The intuition is as follows: if a market participant is willing to incur large slippage losses by submitting a large buy order, they likely want to build their position quickly, potentially in anticipation of the pump. Based on this idea, we compute slippage for each aggregated trade using the formula in \eqref{eq:slippage_loss} below:

\begin{equation}
    \begin{aligned}
        L &= \left|\sum_{i=1}^{N} q_i P_i-P_0\sum_{i=1}^{N}q_i\right|, \\
        L_{\mathrm{signed}} &= L\,\operatorname{sign}\!\left(\sum_{i=1}^{N}s_iq_i\right).
    \end{aligned}
    \label{eq:slippage_loss}
\end{equation}

where $q_i$, $P_i$, and $s_i$ are respectively the unsigned base quantity, execution price, and taker-side sign of the $i^{th}$ print in an aggregated timestamp group, and $P_0$ is its first execution price. Thus $L$ is the absolute difference between observed quote notional and the notional implied by filling the whole group at its first price; the sign is attached from net buyer- or seller-initiated base quantity. The slippage-imbalance feature is $\sum L_{\mathrm{signed}}/\sum L$. As an order consumes liquidity it can appear as a sequence of prints at worsening prices, so this statistic is intended to proxy how aggressively the group walks available liquidity.

\subsubsection{Arithmetic return features}

For each horizon $x$, the implementation computes the arithmetic return $10^4(p_{\mathrm{last}}/p_{\mathrm{prev}}-1)$ over $[T-15\text{min}-x,T-15\text{min})$. Return z-scores compare that window with the mean and sample standard deviation of decision-anchored hourly returns over the preceding 30 days. For $x\in\{5,15\}$ minutes, the exact-window return is converted to an hourly rate before standardization; for horizons of at least one hour, complete cutoff-anchored hourly bars are averaged. Thus the 5-minute, 15-minute, and 1-hour columns are distinct, while every numerator and normalizer remains strictly pre-decision.

\subsubsection{Power law features. Quantile spread of order sizes}

To summarize the tail of aggregated trade notionals $x\geq x_{\min}$ within each lookback window, we use the Pareto density

\begin{equation}
    f(x;\alpha,x_{\min})=(\alpha-1)x_{\min}^{\alpha-1}x^{-\alpha}, \qquad \alpha>1,
    \label{eq:powerlaw_pdf}
\end{equation}

where $x_{\min}$ is the minimum positive aggregated notional observed in the window. The code uses the closed-form maximum-likelihood estimate

\begin{equation}
    \hat{\alpha}=1+\frac{n}{\sum_{i=1}^{n}\log(x_i/x_{\min})}.
    \label{eq:powerlaw_ols}
\end{equation}

Smaller $\hat{\alpha}$ indicates a heavier notional tail. Before cross-sectional standardization, the implementation clips $\hat{\alpha}$ to the closed interval $[1,2]$. This deterministic bound handles degenerate short windows in which the denominator is zero (and the raw estimate is infinite) and limits the influence of extremely sparse tails. This is a fixed model-preprocessing rule, not a sample-dependent 95\% winsorization; no histogram-density regression or log--log OLS is performed.
\\ \\
All of the features discussed are further explained below (see Table~\ref{table:feature_descriptions}):

\begin{table*}[!t]
    \centering
    \small
    \renewcommand{\arraystretch}{1.25}
    \setlength{\tabcolsep}{10pt}
    \begin{tabular}{p{3cm}p{7cm}l}
        \hline
        \textbf{Feature} & \textbf{Description} & \textbf{Notation} \\
        \hline
        Asset return & return of the asset calculated over the window of size (T-15min-x, T-15min) prior to the pump announcement& \textit{asset\_return}[x]@\textit{t} \\
        \hline
        Standardized return & exact-window hourly-rate return (5m, 15m) or mean cutoff-anchored hourly return ($x\geq1$h), minus the 30-day hourly mean and divided by its standard deviation & \textit{asset\_return\_zscore}[x]@\textit{t} \\
        \hline
        Standardized volume in BTC & exact-window hourly-rate volume (5m, 15m) or mean cutoff-anchored hourly volume ($x\geq1$h), standardized by 30-day hourly moments & \textit{quote\_abs\_zscore}[x]@\textit{h} \\
        \hline
        Share of long trades & share of long trades as the ratio of number of long trades and all trades of this time window (T-15min-x, T-15min) & \textit{share\_of\_long\_trades}[x]@\textit{h} \\
        \hline
        Number of trades & number of timestamp-aggregated trades within $[T-15\mathrm{min}-x,T-15\mathrm{min})$ & \textit{num\_trades}[x]@\textit{h} \\
        \hline
        Flow imbalance ratio & volume imbalance ratio computed over the window of size (T-15min-x, T-15min) & \textit{flow\_imbalance}[x]@\textit{h} \\
        \hline
        Quote slippage imbalance ratio & imbalance ratio of slippages in BTC calculated over the same time window & \textit{slippage\_imbalance}[x]@\textit{t} \\
        \hline
        Powerlaw alpha & estimated alpha from powerlaw probability density function using the window of size (T-15min-x, T-15min) prior to the pump & \textit{powerlaw\_alpha}[x]@\textit{h} \\
        \hline
        Number of previous pumps & lifetime number of manifest events for this token strictly before the current announcement & \textit{num\_prev\_pump} \\
        \hline
    \end{tabular}
    \caption{These are the features included in all models. These features are used to describe historical dynamics of price and trading volume in the run up to the pump. $x \in \{5m, 15m, 1h, 2h, 4h, 12h, 1d, 2d, 7d, 14d\}$}
    \label{table:feature_descriptions}
\end{table*}

\subsection{Cross-sectional approach}
\label{sec:cs_standardization}

Given that our event dataset spans early 2018 through August 2022, we account for temporal heterogeneity by applying cross-sectional standardization. This approach reduces market-regime level shifts and focuses the models on relative cross-asset dynamics.

We apply cross-sectional standardization to the unbounded return, return-z-score, quote-volume-z-score, and power-law features. For each selected feature, we subtract its mean within an event and divide by its within-event sample standard deviation:

\begin{equation}
    X_{\text{standardized}} = \frac{X - \bar{X}_{CS}}{\bar{\sigma}_{CS}}
    \label{eq:cs_standardization}
\end{equation}

If a selected feature is constant within an event, it contains no relative ranking information and is mapped to zero; all-missing event columns remain missing for native handling or train-only imputation. We do not apply this second-stage transform to bounded imbalance/share features or to count features (trade counts and lifetime prior-pump count), whose nonnegative levels retain their direct interpretation. Fig.~\ref{fig:data_with_crossnorm} demonstrates the transform on a sample of five cross-sections for three selected numerical features:

\begin{figure*}[!t]
    \centering
    \includegraphics[scale=0.4]{images/cs_standardization_impact.png}
    \caption{Features with and without cross-sectional standardization: histograms and kernel-density estimates for five P\&D events. Applying standardization helps to align distributions more closely with one another across pumps.}
    \label{fig:data_with_crossnorm}
\end{figure*}

\section{Modelling}

In this section, we attempt to create a model able to label this fraudulent activity $15$ minutes before the announcement. As discussed earlier this is a challenging task due to high imbalance in the data. The task of finding manipulated assets can be designed as a binary classification problem where ``1'' implies that the asset is pumped and ``0'' that it is not. Therefore, we trained logistic regression, random forest, and CatBoost classifiers. The problem can also be formulated as learning to rank: the CatBoost ranker is supervised by each asset's realized arithmetic return over the first five minutes after announcement, ranked highest-first within its event cross-section (ordinal position 1 is the highest realized return). These post-announcement returns are labels only; the ranker's regressors are the same strictly pre-decision microstructure features used by the classifiers.

\subsection{Evaluation metrics}
\label{sec:evaluation_metrics}
Top@k is the probability that the ground-truth target occurs among a model's $k$ highest scores. In portfolio terms, it is the probability that a size-$k$ portfolio contains the pump target. Top@k\% replaces the fixed count with a percentage of each cross-section, controlling for the growth in Binance listings described in Section~\ref{section:pump_and_dump_data}. We use one notation convention throughout: Top@k and Top@k\% denote the fixed-count and percentage hit-rate curves; Top@kAUC and Top@k\%AUC denote their respective areas. A concrete threshold replaces $k$ in a point label, so Top@10 is the top-10 hit rate and Top@5\% is the top-5\% hit rate. The reported and optimized area in this study is Top@k\%AUC; Top@kAUC is reserved for an area over fixed-count $k$ and is not reported here.

The area under the Top@k\% curve (Top@k\%AUC) serves as our primary optimization target. In our setup we integrate the Top@k\% curve over the low-$k\%$ region $k\% \in (0,\,20\%]$ and normalize by the length of the integration range so the metric remains in $(0,1)$:
\begin{equation}
    \mathrm{Top@k\%AUC} \;=\; \frac{1}{0.20}\int_{0}^{0.20} \mathrm{Top@k\%}\,d(k\%).
    \label{eq:topkpauc}
\end{equation}
Restricting integration to the $(0,20\%]$ region is motivated by two considerations. First, it is the economically relevant region for a concentrated P\&D portfolio: a practical strategy holds at most the top $10$--$20\%$ of the cross-section in any given event. Second, all reasonably-trained models saturate above $k\% \approx 30\%$ and their Top@k\% curves converge, so integrating over the full $(0,1)$ range dilutes the metric with a large $k$ region where every model is close to $1$ and no longer discriminates between them. Truncating at $20\%$ makes Top@k\%AUC materially more sensitive to differences between models and to hyperparameter choices, which is the property we want from both an Optuna tuning target and an early-stopping signal inside CatBoost. Numerically, Top@k\%AUC values reported in this paper are therefore not directly comparable to AUCs integrated over $(0,1)$; a model that saturates quickly below $20\%$ can attain a Top@k\%AUC close to $0.8$ under this convention.

\subsubsection{Limitations of Top@k-based metrics}
While Top@k metrics are directly relevant for portfolio construction (how many assets must we hold to catch the pump?), they have limitations. They do not measure overall classification quality across the full score distribution and can mask poor calibration. A model may achieve high Top@k by ranking the pump target above most non-targets without producing well-separated probabilities. To address this, we complement Top@k with standard classification metrics.

\subsubsection{Standard classification metrics}
We additionally report three standard metrics on both the validation and test sets:
\begin{itemize}
    \item \textbf{PR-AUC:} Area under the precision--recall curve, which is more informative than receiver operating characteristic (ROC) AUC under extreme class imbalance~\cite{elkan2001cost}.
    \item \textbf{F1-score:} Harmonic mean of precision and recall, computed using a ``Top@1 per cross-section'' decision rule (predicting the highest-scored asset as positive).
    \item \textbf{Balanced accuracy:} Average of sensitivity and specificity, providing a class-imbalance-aware accuracy measure.
\end{itemize}

\subsection{Baseline model}

We use logistic regression as a simple baseline and compare it with a random-ranking model that selects candidate targets uniformly at random within each event cross-section. This comparison tests whether the engineered features contain predictive information beyond random selection. The only baseline adjustment was to set \textit{class\_weight} for the minority class, penalizing false negatives more heavily. We trained the model using the entire training set and were able to obtain sufficient prediction quality, presented in Table~\ref{table:results_topk}. Our primary evaluation metrics are Top@k with values of $k \in \{1, 2, 3, 5, 10, 20, 30\}$ and Top@k\%AUC. These metrics are analyzed in Section \ref{section:results}.

\subsection{Training}
We used the training set to train the models. To reduce overfitting, tree depth was limited through the \textit{max\_depth} hyperparameter for both random forest and CatBoost models. We also applied early stopping for both CatBoost classifier and CatBoost ranker and used validation set to compute evaluation metrics. We created a custom evaluation metric Top@k\%AUC, this metric was used by the trainer to find the optimal number of boosted trees. We also set the number of early stopping rounds to 50 which implies that there is no improvement in Top@k\%AUC for more than 50 rounds, the training process will stop and the model will rollback to the best state. This was used in the training process of the model CatBoost classifier + Top@k\%AUC early stopping (early stopping model with custom evaluation metric).

For each model, hyperparameters were optimized on a held-out validation set, while the test set was reserved exclusively for final evaluation. Then, we used the \textit{Optuna} package \cite{optuna_2019} which implements Bayesian optimization and enables efficient exploration of hyperparameter configurations. Each pipeline was trained $100$ times and evaluated on the validation sample, the set of parameters producing the best results in terms of the chosen Top@k\%AUC metric was chosen and used to train the model. After fitting every candidate, we performed one final model-selection step using validation Top@k\%AUC across all tuned and untuned variants. We froze the selected model before computing any test metric or portfolio result; no test outcome was used for hyperparameter tuning, early stopping, model selection, or portfolio-model selection.

We also used SMOTE proposed by \cite{chawla2002smote} to tackle the problem of class imbalance. It allows to create new observations for the minority class to balance the dataset. Instead of copying existing minority points, it generates new, slightly different ones by interpolating between neighbors.


\section{Results}
\label{section:results}

\subsection{Validation Performance and Model Selection}

Table~\ref{table:validation_model_selection} reports the complete validation comparison used for final model selection. CatBoost + ES achieves the highest validation Top@k\%AUC ($0.667$), ahead of tuned CatBoost ($0.655$) and tuned random forest ($0.628$), and is therefore frozen as the model used for the test portfolio. Its validation Top@1, Top@5, and Top@10 values are $0.175$, $0.427$, and $0.583$, respectively; validation Top@1\%, Top@5\%, and Top@10\% are $0.320$, $0.641$, and $0.718$. It also leads the validation classification metrics with PR-AUC $0.098$, F1 $0.175$, and balanced accuracy $0.586$. These results are model-selection evidence, whereas the test results in the following subsections are final holdout estimates.

\begin{table*}[!t]
    \centering
    \scriptsize
    \setlength{\tabcolsep}{3pt}
    \renewcommand{\arraystretch}{1.15}
    \caption{Validation metrics and final model selection. Top@k\%AUC is the selection criterion; the maximum in each metric column is shown in bold. F1 and balanced accuracy use the Top@1-per-cross-section decision rule.}
    \resizebox{\textwidth}{!}{%
    \begin{tabular}{lrrrrrrrrrr}
        \toprule
        & Top@k\%AUC & Top@1 & Top@5 & Top@10 & Top@1\% & Top@5\% & Top@10\% & PR-AUC & F1 & Balanced accuracy \\
        \midrule
        CatBoost classifier + Top@k\%AUC early stopping & \bfseries 0.667 & \bfseries 0.175 & \bfseries 0.427 & \bfseries 0.583 & 0.320 & \bfseries 0.641 & \bfseries 0.718 & \bfseries 0.098 & \bfseries 0.175 & \bfseries 0.586 \\
        CatBoost classifier + Tuned & 0.655 & 0.165 & 0.417 & 0.524 & 0.301 & 0.592 & 0.709 & 0.087 & 0.165 & 0.581 \\
        Random forest + Tuned & 0.628 & \bfseries 0.175 & 0.398 & 0.485 & \bfseries 0.330 & 0.534 & 0.689 & 0.085 & \bfseries 0.175 & \bfseries 0.586 \\
        Logistic regression + Tuned & 0.616 & 0.107 & 0.301 & 0.485 & 0.223 & 0.563 & 0.680 & 0.059 & 0.107 & 0.552 \\
        CatBoost classifier & 0.615 & 0.165 & 0.350 & 0.505 & 0.291 & 0.544 & 0.670 & 0.080 & 0.165 & 0.581 \\
        Logistic regression & 0.598 & 0.136 & 0.340 & 0.476 & 0.243 & 0.495 & 0.660 & 0.074 & 0.136 & 0.566 \\
        CatBoost ranker + Tuned & 0.550 & 0.107 & 0.301 & 0.466 & 0.223 & 0.476 & 0.602 & 0.048 & 0.107 & 0.552 \\
        Random forest & 0.524 & 0.165 & 0.282 & 0.398 & 0.243 & 0.427 & 0.563 & 0.052 & 0.165 & 0.581 \\
        CatBoost classifier + SMOTE + Tuned & 0.524 & 0.107 & 0.320 & 0.437 & 0.233 & 0.456 & 0.544 & 0.060 & 0.107 & 0.552 \\
        CatBoost ranker & 0.498 & 0.087 & 0.359 & 0.437 & 0.262 & 0.466 & 0.524 & 0.057 & 0.087 & 0.542 \\
        CatBoost classifier + SMOTE & 0.382 & 0.078 & 0.165 & 0.262 & 0.136 & 0.272 & 0.417 & 0.030 & 0.078 & 0.537 \\
        \bottomrule
    \end{tabular}%
    }
    \label{table:validation_model_selection}
\end{table*}

Fig.~\ref{fig:validation_topk} shows validation Top@k against the validation-specific random-ranking baseline, while Fig.~\ref{fig:validation_topkp} shows validation Top@k\% curves and their Top@k\%AUC values. Fig.~\ref{fig:validation_classification} provides the corresponding validation classification comparison. Together with Table~\ref{table:validation_model_selection}, these figures make the selection decision auditable without consulting the test set.

\begin{figure*}[!t]
    \centering
    \includegraphics[scale=0.5]{images/topk_accuracy_validation.png}
    \caption{Top@k accuracy on the validation sample for each tuned baseline, CatBoost + ES, and the validation random-ranking baseline. This figure is computed before the test split is opened.}
    \label{fig:validation_topk}
\end{figure*}

\begin{figure*}[!t]
    \centering
    \includegraphics[scale=0.5]{images/topkp_auc_validation.png}
    \caption{Top@k\% accuracy and Top@k\%AUC on the validation sample. CatBoost + ES has the largest validation Top@k\%AUC ($0.667$) and is selected for the final test portfolio.}
    \label{fig:validation_topkp}
\end{figure*}

\begin{figure*}[!t]
    \centering
    \includegraphics[scale=0.5]{images/classification_metrics_validation.png}
    \caption{PR-AUC, F1, and balanced accuracy on the validation sample. F1 and balanced accuracy use the Top@1-per-cross-section decision rule.}
    \label{fig:validation_classification}
\end{figure*}

\subsection{Top@k accuracy}
After freezing CatBoost + ES as the portfolio model on validation, we opened the test split and computed final Top@k accuracy for every candidate. Test comparisons among non-selected models are descriptive and do not alter the selection decision. %The results are presented in Table~\ref{table:results_topk}.

\begin{table*}[!t]
    \centering
    \setlength{\tabcolsep}{6pt} % Adjust the column spacing
    \renewcommand{\arraystretch}{1.2} % Adjust the row spacing
    \caption{Top@k accuracy of models on test sample}
    \begin{tabular}{lcccccc}
        \toprule
        Model & Top@1 & Top@2 & Top@5 & Top@10 & Top@20 & Top@30 \\
        \midrule
        Logistic regression & 0.113 & 0.150 & 0.388 & 0.550 & 0.663 & 0.700 \\
        Logistic regression + Tuned & 0.100 & 0.200 & \textbf{0.425} & \textbf{0.575} & \textbf{0.700} & \textbf{0.763} \\
        Random forest & 0.075 & 0.125 & 0.275 & 0.363 & 0.575 & 0.675 \\
        Random forest + Tuned & 0.113 & 0.175 & 0.350 & 0.500 & 0.650 & 0.713 \\
        CatBoost classifier & 0.113 & 0.175 & 0.388 & 0.463 & 0.575 & 0.688 \\
        CatBoost classifier + Tuned & 0.138 & \textbf{0.275} & 0.400 & 0.525 & 0.650 & 0.725 \\
        CatBoost classifier + SMOTE & 0.063 & 0.088 & 0.150 & 0.213 & 0.313 & 0.400 \\
        CatBoost classifier + SMOTE + Tuned & 0.113 & 0.163 & 0.300 & 0.425 & 0.538 & 0.600 \\
        CatBoost ranker & 0.100 & 0.138 & 0.288 & 0.363 & 0.538 & 0.625 \\
        CatBoost ranker + Tuned & 0.063 & 0.163 & 0.275 & 0.463 & 0.613 & 0.700 \\
        CatBoost classifier + Top@k\%AUC early stopping & \textbf{0.150} & 0.225 & 0.388 & 0.563 & 0.663 & 0.750 \\
        \bottomrule
    \end{tabular}
    \label{table:results_topk}
\end{table*}


Table~\ref{table:results_topk} shows that logistic regression remains a competitive baseline, indicating that the engineered features contain predictive signal. CatBoost + ES leads at $k=1$ ($0.150$); tuned CatBoost leads at $k=2$ ($0.275$); and tuned logistic regression leads at $k=5$ ($0.425$), $k=10$ ($0.575$), $k=20$ ($0.700$), and $k=30$ ($0.763$). SMOTE-based variants remain below the leading class-weighted models across the reported thresholds, suggesting that synthetic oversampling is poorly aligned with the cross-sectional microstructure features used in this setting. Fig.~\ref{fig:results_topk} reports the corresponding Top@k curves.



\begin{figure*}[!t]
    \centering
    \includegraphics[scale=0.5]{images/topk_accuracy_test.png}
    \caption{Top@k accuracy on the test sample for each tuned baseline, CatBoost + ES, and the random-ranking baseline. The horizontal axis is the portfolio size $k$; the vertical axis is the empirical probability that the manipulation target appears among the $k$ highest-scored assets. Table~\ref{table:results_topk} reports both tuned and untuned variants.}
    \label{fig:results_topk}
\end{figure*}

All trained models outperform the random-ranking baseline, indicating that the feature set contains predictive information about future manipulation targets. Top@k values for the random-ranking model are computed as follows:

\begin{equation}
    \mathrm{Top@k}_{\text{random}} = \frac{1}{N} \sum_{i=1}^{N}{\frac{k}{T_{i}}}
\end{equation}

where $N$ is the number of cross-sections in the test set and $T_i$ is the $i$-th cross-section size. This implies that our models are able to consistently capture underlying patterns in the data that describe P\&D events. After 100 tuning trials, the ranking model is competitive but remains below the leading classification models on the aggregate ranking metric (test Top@k\%AUC $0.629$); the mechanism is analyzed later in this section.

\subsection{Top@k\% accuracy}

Table~\ref{table:results_topkp} reports the corresponding Top@k\% pattern. Tuned CatBoost leads at $k\%=1\%$ ($0.338$). Untuned logistic regression leads at $k\%=2\%$ ($0.488$), while tuned logistic regression leads at $k\%=5\%$ ($0.638$), $10\%$ ($0.750$), and $20\%$ ($0.863$). At $k\%=50\%$, tuned logistic regression and untuned random forest tie at $0.963$. The aggregate ranking differences are assessed in Section~\ref{sec:significance}.

%Then, we estimated Top@k\% for all models using the same test sample, the results are reported in Table~\ref{table:results_topkp}:

\begin{table*}[!t]
    \centering
    \setlength{\tabcolsep}{6pt}
    \renewcommand{\arraystretch}{1.2}
    \caption{Top@k\% accuracy of models on test sample}
    \begin{tabular}{lcccccc}
        \toprule
        Model & Top@1\% & Top@2\% & Top@5\% & Top@10\% & Top@20\% & Top@50\% \\
        \midrule
        Logistic regression & 0.275 & \textbf{0.488} & 0.613 & 0.700 & 0.838 & 0.950 \\
        Logistic regression + Tuned & 0.325 & 0.475 & \textbf{0.638} & \textbf{0.750} & \textbf{0.863} & \textbf{0.963} \\
        Random forest & 0.188 & 0.325 & 0.488 & 0.675 & 0.763 & \textbf{0.963} \\
        Random forest + Tuned & 0.275 & 0.425 & 0.575 & 0.700 & 0.825 & 0.925 \\
        CatBoost classifier & 0.325 & 0.413 & 0.513 & 0.675 & 0.800 & 0.950 \\
        CatBoost classifier + Tuned & \textbf{0.338} & 0.425 & 0.625 & 0.725 & 0.813 & 0.950 \\
        CatBoost classifier + SMOTE & 0.113 & 0.175 & 0.288 & 0.388 & 0.563 & 0.800 \\
        CatBoost classifier + SMOTE + Tuned & 0.238 & 0.325 & 0.513 & 0.588 & 0.763 & 0.900 \\
        CatBoost ranker & 0.200 & 0.338 & 0.450 & 0.613 & 0.725 & 0.925 \\
        CatBoost ranker + Tuned & 0.200 & 0.375 & 0.525 & 0.700 & 0.825 & 0.888 \\
        CatBoost classifier + Top@k\%AUC early stopping & 0.313 & 0.463 & 0.600 & 0.738 & 0.838 & 0.950 \\
        \bottomrule
    \end{tabular}
    \label{table:results_topkp}
\end{table*}


\subsection{Top@k\%AUC}



In order to control for the problem of various cross-section sizes, we  want to compute Top@k\% instead of regular Top@k. Taking the normalized area under the Top@k\% curve over $k\% \in (0, 20\%]$ allows us to compare models in the economically relevant low-portfolio-size region, consistent with Eq.~\ref{eq:topkpauc}. In Fig.~\ref{fig:results_topkp} we compare the models based on this metric.

\begin{figure*}[!t]
    \centering
    \includegraphics[scale=0.5]{images/topkp_auc_test.png}
    \caption{Top@k\% accuracy on the test sample, normalizing the portfolio size by the cross-section size to control for its growth over time. The horizontal axis is the fraction $k\%$ of the cross-section retained; the vertical axis is the probability that the target is covered.}
    \label{fig:results_topkp}
\end{figure*}

\subsection{Classification Metrics}

We complement the ranking-based Top@k metrics with three standard classification metrics (defined in Section~\ref{sec:evaluation_metrics}): PR-AUC, F1-score (Top@1-per-cross-section decision rule), and balanced accuracy. PR-AUC is preferred over ROC-AUC at this imbalance ratio~\cite{elkan2001cost}. Fig.~\ref{fig:classification_metrics} reports all three for the tuned baselines and CatBoost + ES; the underlying comparison also evaluates the untuned variants. Across all variants, tuned CatBoost leads PR-AUC ($0.0866$), followed by CatBoost + ES ($0.0836$). CatBoost + ES leads F1 ($0.1500$) and balanced accuracy ($0.5735$), followed by tuned CatBoost at $0.1375$ and $0.5672$. This pattern is consistent with the strong low-$k$ ranking performance of both CatBoost variants.

\begin{figure*}[!t]
    \centering
    \includegraphics[scale=0.5]{images/classification_metrics_test.png}
    \caption{Standard classification metrics on the test set for all tuned models and CatBoost + ES. Each grouped bar reports Precision-Recall AUC (left), F1-score under a Top@1-per-cross-section decision rule (middle), and balanced accuracy, the average of sensitivity and specificity (right). PR-AUC is preferred over ROC-AUC under the $\sim 1{:}287$ class imbalance.}
    \label{fig:classification_metrics}
\end{figure*}

\subsection{Analysis of SMOTE Performance}

SMOTE-based CatBoost variants performed worse than the leading class-weighted classifiers. The untuned and tuned SMOTE variants reach $6.25\%$ and $11.25\%$ Top@1, with Top@k\%AUC of $0.378$ and $0.574$, respectively; tuned logistic regression reaches the strongest test Top@k\%AUC of $0.697$. Thus, 100-trial tuning materially improves the SMOTE variant but does not close the gap. We attribute the remaining deficit to a structural mismatch between SMOTE's generative assumption (local linear interpolation between minority samples in feature space) and three specific properties of our panel microstructure features.

\begin{enumerate}
    \item \textbf{Bounded and sign-meaningful features.} Several of our features---flow imbalance, slippage imbalance, share of buy-initiated trades---are naturally bounded in $[-1,1]$ (or $[0,1]$) and are informative primarily near the positive extreme (e.g., heavy buy-side imbalance just before a pump). Interpolating between two pumped samples whose imbalance is $+0.85$ and another whose imbalance is $-0.20$ yields a synthetic ``pump'' with imbalance $+0.32$, which lies inside the bulk of the negative-class distribution. SMOTE therefore dilutes rather than reinforces the class-discriminative region of the feature space.
    \item \textbf{Cross-sectional reference.} The unbounded return, standardized-volume, and tail-shape features are standardized within the event cross-section (Eq.~\ref{eq:cs_standardization}); for those columns, a value of $+2$ means ``two cross-sectional standard deviations above the same-event average.'' Bounded ratios and count features retain their direct scales but are still interpreted jointly with those relative columns. SMOTE's nearest-neighbor interpolation mixes complete feature vectors from \emph{different} events, producing synthetic samples whose relative and absolute components have no single corresponding cross-section in which they are coherent. The synthetic samples can be plausible marginally but out-of-distribution jointly with respect to the panel structure.
    \item \textbf{High-dimensional sparsity.} With 81 features and only 290 positive training samples, the nearest-neighbor graph that SMOTE relies on is sparse and noisy: each minority point has few true neighbors in its class, and interpolation typically bridges across different sub-populations (e.g., small-cap vs.\ medium-cap pumps). This is the regime in which Blagus and Lusa~\cite{blagus2013smote} showed that SMOTE is at best neutral and frequently \emph{degrades} classification performance.
\end{enumerate}

By contrast, class-weight reweighting in the CatBoost loss function does not add, remove, or perturb any feature value. It only changes the per-sample gradient magnitude, concentrating the optimizer on the decision boundary between the real positive and negative samples in their native, cross-sectionally-normalized feature space. This explains the consistently superior performance of class-weight-based imbalance handling over synthetic oversampling in our setting and is consistent with the cost-sensitive-learning literature~\cite{elkan2001cost}.

\subsection{Analysis of Ranking Model Performance}

The 100-trial-tuned CatBoost ranker is competitive but does not lead on the aggregate ranking metric. Its highest-first return-rank target yields test Top@k\%AUC of $0.629$ and Top@1 of $6.25\%$, while the untuned variant attains $0.548$ and $10.0\%$, respectively. Tuning therefore closes much of the aggregate gap, although both variants remain below CatBoost + ES ($0.674$) and tuned logistic regression ($0.697$). A natural expectation from the learning-to-rank literature~\cite{poh_ltr}, where listwise ranking outperforms pointwise regression in cross-sectional equity alpha prediction, is that ranking should dominate here as well. That expectation does not hold on the aggregate ranking metric in our setting, and the deviation is informative.

\textbf{Target-signal density.}
In the cross-sectional LTR setup of~\cite{poh_ltr}, every asset in the cross-section carries an informative continuous return target, so the listwise loss receives dense, ordinally-meaningful supervision. In our setup, exactly one asset per cross-section is a true manipulation target. All other assets contribute noise to the ordering: their 5-minute post-announcement returns are driven by background market movements, independent news, or incidental volatility, none of which the model is asked to predict. The ranker therefore spends most of its gradient on nuisance ordering among non-pumped assets, while the classifier only has to separate the single pump from the rest.

\textbf{Loss-function alignment with the Top@k evaluation objective.}
Pairwise and listwise ranking losses (YetiRank, PairLogit, etc.) integrate pairwise or position-weighted comparisons across the entire list. When the positive cardinality is 1, most of those pairs are negative-vs.-negative and therefore non-informative for Top@k. The classification log-loss, conversely, concentrates gradient mass on the decision boundary between the single positive and the $\sim$290 negatives, which is exactly the boundary that Top@k accuracy measures. Class-weight reweighting further tilts the log-loss toward the positive class, yielding a gradient that is sample-efficient in the 1-vs.-many regime.

\textbf{Target-construction noise.}
The continuous ranking target (realized arithmetic return over the first five minutes after announcement) is additionally noisy because (i) coincident market-wide moves can push a non-pumped asset's five-minute return above the pump target's, and (ii) micro-cap order books produce large relative returns on very small trades, which the ranker has no way to discount. In the classification formulation the binary label is immune to this type of noise.

Taken together, these three factors explain why, in the extreme-imbalance cross-sectional setting that characterizes P\&D target prediction, classification with class-weight reweighting dominates listwise ranking even though the evaluation metric is itself a ranking metric.

In terms of the aggregated ranking metric, tuned logistic regression attains the highest test Top@k\%AUC point estimate ($0.697$), followed by CatBoost + ES ($0.674$), tuned CatBoost ($0.662$), untuned logistic regression ($0.660$), and tuned random forest ($0.651$). For the pre-specified paired comparison, the non-ES comparator is chosen exclusively by validation Top@k\%AUC; this selects tuned CatBoost ($0.655$ on validation), not the test-leading logistic model. Section~\ref{sec:significance} reports that paired test. At the most selective percentage threshold tuned CatBoost leads Top@1\% ($0.338$); untuned logistic regression leads Top@2\% ($0.488$); and tuned logistic regression leads Top@5\% ($0.638$), Top@10\% ($0.750$), and Top@20\% ($0.863$) (Table~\ref{table:results_topkp}).

\subsection{Interpretation of the models}

This section shifts from model comparison to model interpretation by analyzing which features contribute most strongly to the CatBoost + ES predictions. We use CatBoost PredictionValuesChange importance. In the full-training fit shown in Fig.~\ref{fig:catboost_fi}, \textit{num\_prev\_pump} is the leading feature, followed by \textit{num\_trades@7D}, \textit{num\_trades@14D}, \textit{asset\_return\_zscore@2D}, and \textit{share\_of\_long\_trades@1H}. The dominance of prior-pump-history and trade-count variables is consistent with manipulators repeatedly selecting familiar, liquid-enough tokens. The multi-day standardized return and one-hour aggressor-side share provide complementary evidence of unusual pre-announcement activity.

To verify that this is not an artifact of a single fit, we recomputed the feature importances inside each of the 25 random 70\% training-subset retraining runs of Section~\ref{sec:robustness}, which yields a training-subset distribution of importance for every feature. Table~\ref{table:feature_importance_quantiles} reports the per-feature percentiles of the 15 most important features, sorted by median. Across subsets, \textit{num\_prev\_pump} has the highest median importance, followed by \textit{num\_trades@14D}, \textit{num\_trades@7D}, \textit{num\_trades@2D}, and \textit{asset\_return@4H}; thus the broad prior-history and activity pattern is stable even though the precise ordering changes across fits.

\begin{figure*}[!t]
    \centering
    \includegraphics[scale=0.5]{images/feature_importances.png}
    \caption{Feature importances for the CatBoost classifier + Top@k\%AUC early stopping. Bars report CatBoost PredictionValuesChange for the full-training fit; the leading features are prior pump history, 7-day and 14-day trade counts, two-day standardized return, and the one-hour share of long trades.}
    \label{fig:catboost_fi}
\end{figure*}

\begin{table*}[!t]
    \centering
    \setlength{\tabcolsep}{6pt}
    \renewcommand{\arraystretch}{1.2}
    \caption{Training-subset distribution of CatBoost + ES feature importances (PredictionValuesChange) across the 25 random 70\% training-subset retraining runs of Section~\ref{sec:robustness}. The 15 features with the highest median importance are shown; columns report per-feature percentiles and rows are sorted by median.}
    \begin{tabular}{lrrrrr}
        \toprule
        & 10\% & 25\% & Median & 75\% & 90\% \\
        Feature & & & & & \\
        \midrule
        num\_prev\_pump & 8.74 & 12.58 & 16.28 & 20.57 & 25.88 \\
        num\_trades@14D & 2.69 & 4.03 & 7.11 & 8.68 & 13.04 \\
        num\_trades@7D & 2.55 & 3.72 & 5.92 & 8.67 & 10.91 \\
        num\_trades@2D & 0.50 & 2.17 & 5.02 & 8.57 & 16.60 \\
        asset\_return@4H & 0.45 & 1.84 & 2.68 & 4.12 & 4.82 \\
        slippage\_imbalance@12H & 0.00 & 0.83 & 2.60 & 4.87 & 8.76 \\
        slippage\_imbalance@2D & 0.44 & 1.12 & 2.35 & 3.83 & 7.52 \\
        asset\_return@2H & 0.66 & 1.80 & 2.25 & 4.08 & 5.37 \\
        quote\_abs\_zscore@1H & 0.26 & 1.14 & 2.24 & 2.95 & 4.49 \\
        asset\_return@12H & 0.19 & 1.22 & 2.06 & 2.62 & 4.65 \\
        num\_trades@1D & 0.00 & 0.61 & 1.83 & 2.92 & 4.02 \\
        asset\_return\_zscore@12H & 0.00 & 0.48 & 1.73 & 2.63 & 4.36 \\
        asset\_return\_zscore@1H & 0.06 & 0.78 & 1.69 & 2.70 & 4.09 \\
        asset\_return\_zscore@4H & 0.03 & 0.57 & 1.51 & 2.85 & 3.53 \\
        asset\_return@1D & 0.00 & 0.66 & 1.48 & 1.95 & 3.04 \\
        \bottomrule
    \end{tabular}
    \label{table:feature_importance_quantiles}
\end{table*}

\subsection{Statistical Significance}
\label{sec:significance}

To quantify uncertainty around the observed performance differences, we employ cross-section-level bootstrap resampling. Each bootstrap sample draws 80 test-set cross-sections with replacement, and all metrics are recomputed on the resampled cross-sections. Resampling at the cross-section level (rather than the row level) is the methodologically appropriate choice here: each pump event is the natural unit of observation, and rows belonging to the same event are highly dependent because they share the cross-section-standardized denominator (Eq.~\ref{eq:cs_standardization}). Row-level bootstrapping would produce optimistically narrow intervals.

Table~\ref{table:significance_ci} reports the Top@k\%AUC point estimate together with its 95\% bootstrap confidence interval for CatBoost + ES and tuned CatBoost, the non-ES comparator selected by validation Top@k\%AUC. The intervals span approximately $0.16$--$0.17$, which reflects both the number of positive test events ($n = 80$) and the fact that the $(0, 20\%]$ integration range of Top@k\%AUC (Eq.~\ref{eq:topkpauc}) concentrates the metric on the steep, high-variance region of the Top@k\% curve.

\begin{table}[!t]
    \centering
    \setlength{\tabcolsep}{6pt}
    \renewcommand{\arraystretch}{1.2}
    \caption{Top@k\%AUC on the test set with 95\% cross-section-level bootstrap confidence intervals (1000 iterations). Top@k\%AUC is integrated and normalized over $k\% \in (0, 20\%]$ (Eq.~\ref{eq:topkpauc}).}
    \begin{tabular}{lcc}
        \toprule
        Model & Top@k\%AUC & 95\% CI \\
        \midrule
        \textbf{CatBoost + ES} & \textbf{0.674} & $[\textbf{0.592},\,\textbf{0.755}]$ \\
        CatBoost classifier + Tuned & 0.662 & $[0.575,\,0.747]$ \\
        \bottomrule
    \end{tabular}
    \label{table:significance_ci}
\end{table}

Fig.~\ref{fig:bootstrap_ci} visualizes the 95\% bootstrap confidence intervals of CatBoost + ES at each $k$ and $k\%$. To test whether the point-estimate ordering is distinguishable from sampling noise, we conduct a \emph{paired} bootstrap hypothesis test: on each of 2000 resamples, CatBoost + ES and validation-selected tuned CatBoost are scored on the same resampled cross-sections. The observed Top@k\%AUC difference (ES minus tuned CatBoost) is $0.0116$ with paired two-sided 95\% CI $[-0.0120,\,0.0369]$ and a one-sided $p$-value of $0.1739$ under the pre-specified alternative that CatBoost + ES is better. We therefore do not reject the null at the 5\% level: CatBoost + ES has the larger point estimate, but its aggregate-ranking advantage over the strongest validation-selected non-ES model is not statistically distinguishable from sampling noise. CatBoost + ES leads F1 and balanced accuracy, while tuned CatBoost leads PR-AUC (Fig.~\ref{fig:classification_metrics}).

\begin{figure*}[!t]
    \centering
    \includegraphics[scale=0.5]{images/significance_bootstrap_ci.png}
    \caption{Cross-section-level bootstrap 95\% confidence intervals for CatBoost + ES on the test set, based on 1000 resamples of the 80 test cross-sections with replacement. Left: Top@k accuracy with error bars at $k\in\{1,2,3,5,10,20,30\}$. Right: Top@k\% accuracy with error bars at $k\%\in\{1,2,5,10,20\}$. The dashed horizontal segment in each panel is the random-ranking baseline at the same $k$ / $k\%$. \textbf{Takeaway:} the point estimates lie above random, while the interval widths argue for cautious interpretation of narrow differences between models.}
    \label{fig:bootstrap_ci}
\end{figure*}

\subsection{Robustness Analysis}
\label{sec:robustness}

\subsubsection{Training data perturbation}
We assess the stability of CatBoost + ES by retraining it 25 times, each time using a random 70\% subset of the training cross-sections while keeping the validation and test sets fixed. Table~\ref{table:robustness_summary} summarizes the distribution of test-set metrics across runs.

\begin{table}[!t]
    \centering
    \setlength{\tabcolsep}{4pt}
    \renewcommand{\arraystretch}{1.2}
    \caption{Robustness analysis: distribution of test-set metrics across 25 runs with random 70\% training subsets. Top@k\%AUC is integrated and normalized over $k\% \in (0, 20\%]$ (Eq.~\ref{eq:topkpauc}). The standard deviation of Top@k\%AUC is $0.013$.}
    \begin{tabular}{lccc}
        \toprule
        Metric & Mean & Std & [10\%, 90\%] \\
        \midrule
        Top@k\%AUC & 0.678 & 0.013 & [0.662, 0.693] \\
        Top@1\% & 0.334 & 0.038 & [0.305, 0.370] \\
        Top@2\% & 0.463 & 0.032 & [0.438, 0.500] \\
        Top@5\% & 0.619 & 0.023 & [0.593, 0.645] \\
        Top@10\% & 0.727 & 0.016 & [0.705, 0.745] \\
        Top@20\% & 0.839 & 0.017 & [0.825, 0.858] \\
        \bottomrule
    \end{tabular}
    \label{table:robustness_summary}
\end{table}

The Top@k\%AUC exhibits a standard deviation of $0.013$ across runs, with the central 80\% of runs falling within $[0.662, 0.693]$. Fig.~\ref{fig:robustness_distribution} visualizes the distribution. This demonstrates that performance is not driven by a single fortunate training-set composition. Because feature ordering and estimator seeds are fixed, the runs and every other reported result in this paper are exactly reproducible from the released code.

\begin{figure*}[!t]
    \centering
    \includegraphics[scale=0.5]{images/robustness_metrics_distribution.png}
    \caption{Distribution of test-set metrics of CatBoost + ES across 25 independent retraining runs. Each run uses a different random 70\% subset of the training cross-sections; the validation and test sets are held fixed. Left: histogram of Top@k\%AUC (integrated and normalized over $k\% \in (0, 20\%]$) across the 25 runs; the vertical dashed line is the across-run mean. Right: box plots of Top@k\% at $k\% \in \{1,2,5,10,20\}$. \textbf{Takeaway:} $\sigma_{\mathrm{Top@k\%AUC}} = 0.013$ and the central 80\% of runs fall in $[0.662,\,0.693]$.}
    \label{fig:robustness_distribution}
\end{figure*}

\subsubsection{Temporal subperiod analysis}
Our test set spans from May~2021 to August~2022, covering the late-2021 bull market and the onset of the 2022 contagion around the Terra/LUNA collapse. We conducted the subperiod analysis with June~1, 2022 chosen \emph{ex ante} as the bull$\rightarrow$bear regime boundary, so that the split is not a function of the model's test-set performance.

The early subperiod $[\text{May 2021},\text{June 1, 2022})$ contains 76 valid pump events and attains Top@k\%AUC of $0.676$. The late subperiod $[\text{June 1, 2022},\text{August 2022}]$ contains only 4 labeled events. We therefore apply a \texttt{min\_pumps=10} guard and explicitly exclude the late subperiod from quantitative comparison: any Top@k\%AUC estimated on four events has a bootstrap 95\% CI too wide to be useful. We acknowledge two implications: (i) performance can be assessed with useful precision only through May~2022 in this sample; and (ii) generalizability to a sustained bear market cannot be concluded and is flagged in Section~\ref{sec:limitations}.

\section{Portfolio strategy}
\label{sec:portfolio}

This section evaluates a portfolio strategy inspired by \cite{pd_anatomy}, with modified construction and execution assumptions tailored to the proposed scoring framework. Authors of \cite{pd_anatomy} proposed using trained classification models to produce logits for each asset in the cross-section, then construct a portfolio of assets that have their logit of being pumped above some threshold. They also proposed using weights of the portfolio proportional to the logits produced by the model. They backtested their strategy using the following assumptions: assets that are not pumped do not change in price over the pump, authors buy at open price of the last hourly candle prior to the pump and are able to sell at half of pump's max return. They backtested their strategy using the test sample from October 29, 2018 to January 11, 2019 and were able to achieve a return of 60\% over approximately two months. We extend this approach by constructing portfolios from the proposed model scores and by adopting more conservative execution assumptions. 

We propose using models tuned for Top@k\%AUC which will ensure that the models are optimized not only for Top@1 accuracy but for any arbitrary values of k. We relaxed some assumptions of \cite{pd_anatomy} in the way we trade our portfolios, the timeline of how we trade each pump is described below:

\begin{enumerate}
  \item Rank tickers by model logits (descending).
  \item Select the Top@k set of tickers.
  \item Submit equal-weight entries 15 minutes before the announcement and fill each at its first subsequent pre-announcement trade; omit a leg if no such trade occurs. The omitted leg's fixed $1/k$ weight remains in cash and is not redistributed.
  \item Hold until the announcement; if not pumped, sell at the announcement; if pumped, sell 1 minute after the pump.
  \item Liquidate the complete acquired base-asset quantity, convert the BTC proceeds to USDT at the exit-time completed indicative rate, and subtract a $25$ bp ($0.25\%$) round-trip cost from entry capital.
\end{enumerate}



We use CatBoost classifier + Top@k\%AUC ES as the primary portfolio model based on validation performance. The execution-aware sensitivity analysis pre-specifies $k=5$; test results for $k\in\{1,2,5,10,20,30\}$ are also reported as descriptive concentration checks and are not used to retune the strategy.


\begin{table*}[h!]
    \centering
    \setlength{\tabcolsep}{5pt}
    \renewcommand{\arraystretch}{1.2}
    \begin{tabular}{lccccc}
        \toprule
        \makecell[c]{Portfolio /\\benchmark} &
        \makecell[c]{Avg. event\\return} &
        \makecell[c]{Cumulative\\return} &
        \makecell[c]{Annualized\\portfolio\\return} &
        \makecell[c]{Annualized\\portfolio\\volatility} &
        \makecell[c]{Sharpe ratio} \\
        \midrule
        $k=1$  & 0.0123 & \textbf{0.9834} & 0.7621 & 0.7423 & 1.0267 \\
        $k=2$  & 0.0118 & 0.9470 & 0.7338 & 0.4161 & 1.7638 \\
        $k=5$  & 0.0120 & 0.9570 & 0.7416 & 0.2896 & \textbf{2.5608} \\
        $k=10$ & 0.0072 & 0.5751 & 0.4457 & 0.2130 & 2.0920 \\
        $k=20$ & 0.0043 & 0.3451 & 0.2674 & 0.1251 & 2.1370 \\
        $k=30$ & 0.0029 & 0.2358 & 0.1828 & 0.0894 & 2.0448 \\
        BTC/USDT buy-and-hold & --- & $-0.5727$ & $-0.4448$ & 0.7340 & $-0.6060$ \\
        \bottomrule
    \end{tabular}
    \caption{Performance by portfolio size and the time-aligned BTC/USDT benchmark}
    \label{table:performance_by_k}
\end{table*}


Table~\ref{table:performance_by_k} and Fig.~\ref{fig:portfolio_test} show that the risk-adjusted return (Sharpe ratio) peaks at the pre-specified $k=5$ portfolio ($2.56$), followed by $k=20$ ($2.14$). The concentrated $k=1$ portfolio has the largest cumulative return ($0.983$) but is much more volatile and has a Sharpe ratio of $1.03$, while broader portfolios dilute the signal. All six portfolio sizes finish with positive no-reinvestment cumulative returns ($0.236$--$0.983$), whereas BTC/USDT loses $57.3\%$ over the same calendar. Portfolio returns use a fixed position size for every event.

\begin{figure*}[!t]
    \centering
    \includegraphics[scale=0.6]{images/portfolios_vs_btc_test.png}
    \caption{Cumulative return paths on the test set for portfolios of size $k \in \{1,2,5,10,20,30\}$, ranked by CatBoost + ES scores, together with a continuous BTC/USDT buy-and-hold benchmark over the identical calendar interval. Portfolio returns are accumulated without reinvestment and include a 25 bps round-trip cost. The benchmark is the solid black line.}
    \label{fig:portfolio_test}
\end{figure*}

For market context, Fig.~\ref{fig:portfolio_test} includes a continuous BTC/USDT buy-and-hold position over the exact strategy calendar, from May~1, 2021 16:45:08 UTC to August~14, 2022 16:01:05 UTC. It enters at the first decision timestamp ($T-15$ minutes) and is valued finally at the last strategy exit ($T+1$ minute), using completed one-minute BTC/USDT VWAPs. Daily UTC observations connect those exact endpoints. Each plotted price increment is expressed relative to the fixed initial BTC price, so its cumulative sum equals the simple full-period BTC return exactly. BTC/USDT records cumulative return $-0.573$, annualized return $-0.445$, daily annualized volatility $0.734$, and Sharpe ratio $-0.606$. This is a calendar-aligned market benchmark rather than an exposure-matched strategy: BTC remains invested between events, whereas each Top@k portfolio is exposed only during its event window.

\subsection{Execution-Aware Backtesting with Price Impact}
\label{sec:price_impact}

The results above use the first observed trade at or after each execution threshold but do not otherwise charge size-dependent slippage, which is unrealistic for the illiquid small-cap tokens typically targeted by P\&D operations. To assess practical feasibility, we extend the backtesting framework with a fitted market-impact model that produces \emph{size-dependent, dynamic} slippage in addition to the fixed round-trip fee. The model serves two purposes: first, it estimates a volume-weighted average price (VWAP) for both entry and exit; second, the shape of the impact curve is itself the \emph{liquidity constraint} of the strategy---larger orders remain executable, but the fitted $\beta\sqrt{Q}$ curve makes each additional unit of notional arrive at progressively worse prices, so the economically relevant question shifts from ``is the order executable?'' to ``at what order size does impact erode the signal?''.

\subsubsection{Impact model specification}
We model price impact in basis points as a function of order notional $Q$ (in USDT):

\begin{equation}
    I(Q) = \beta \sqrt{Q}.
    \label{eq:impact_model}
\end{equation}

The square-root specification is a well-established empirical regularity in market microstructure. Kyle~\cite{kyle1985continuous} introduced the foundational model of linear price impact from informed order flow; subsequent work demonstrated that the aggregate impact of executed orders follows a concave, approximately square-root, relationship with order size. Almgren and Chriss~\cite{almgren2001optimal} formalized this within the optimal execution framework, while Almgren et al.~\cite{almgren2005direct} provided direct empirical estimates of the power-law exponent ($\approx 0.6$) using institutional equity data. Bouchaud et al.~\cite{bouchaud2009markets} surveyed the evidence for the square-root law across asset classes, and T\'{o}th et al.~\cite{toth2011anomalous} demonstrated its universality across stocks and time periods using a large meta-order dataset. In cryptocurrency markets specifically, Donier and Bonart~\cite{donier2015million} confirmed the square-root law in Bitcoin using over one million reconstructed meta-orders, showing that the relationship holds across four decades of order sizes.

We fit a single coefficient $\beta \geq 0$ via constrained OLS (no intercept) for each asset, pooling absolute impacts from both buy and sell candles. The entry regression uses the 14 days of 5-minute candles ending strictly at the actual simulated entry; it therefore cannot observe any later trade. For the manipulated asset, the execution backtest fits 5-second sell-dominated candles from $T$ up to, but excluding, the realized exit timestamp, capped at $T+10$ minutes. In each case, absolute net volume serves as $Q$ and the absolute close-to-close price change (in bps) serves as the impact observation. Entry notionals are converted to USDT using the most recent completed BTC/USDT one-minute bar at entry; the exit model uses the completed rate available at the announcement. Neither conversion falls forward to a future or incomplete bar.

\subsubsection{Impact measurement from candle data}
\textbf{Entry model (5-minute candles).} We aggregate the 14-day trade history ending strictly at the actual simulated entry timestamp into 5-minute candles. For each candle we compute the net buy volume $Q_{\mathrm{net}} = 2 \times V_{\mathrm{taker\,buy}} - V_{\mathrm{total}}$ (in quote currency), where $V_{\mathrm{taker\,buy}}$ is the taker-buy quote volume and $V_{\mathrm{total}}$ is the total quote volume. The absolute net volume $|Q_{\mathrm{net}}|$ serves as the order size, and the absolute close-to-close price change $|p_{\mathrm{close}} / p_{\mathrm{prev\,close}} - 1| \times 10^4$ bps serves as the impact observation. Observations from both buy- and sell-dominated candles are pooled into a single regression, which doubles the effective sample size and produces a more stable estimate of $\beta$. The 5-minute resolution matches the intended execution horizon: we assume the order is worked over a single candle interval, so the candle-level impact directly estimates the cost of execution.

\textbf{Exit model for manipulated assets (5-second candles, sell-only).} For the pumped asset, the backtest fits a separate model using only sell-dominated candles (negative $Q_{\mathrm{net}}$) in $[T,\min(T+10\mathrm{min},t_{\mathrm{exit}}))$. The right-open exit bound prevents post-exit observations from entering the simulated execution cost. The full 10-minute window is used only for the illustrative regression figure below.

We model liquidity through a continuous price-impact penalty rather than a hard participation cap. This makes larger orders executable only at increasingly adverse VWAPs. A hard order-book depth or participation-rate cap would further reduce realized returns and is therefore a conservative extension for future work.

\subsubsection{VWAP entry and exit price estimation}
\label{sec:vwap_estimation}

The execution simulation proceeds as follows for each selected asset:

\textbf{Step 1: Determine raw prices.} Once the model has scored the cross-section at $t_{\mathrm{pump}} - \Delta t_{\mathrm{buy}}$, the entry price $p^{\mathrm{entry}}$ is the first trade price at or after that decision timestamp and strictly before $t_{\mathrm{pump}}$. If no such trade occurs, the leg is not executable and is omitted. The exit price $p^{\mathrm{exit}}$ is the first trade price at or after $t_{\mathrm{pump}}$ (or $t_{\mathrm{pump}} + \Delta t_{\mathrm{sell}}$ for the manipulated asset). Thus no simulated fill predates the signal that generated it.

\textbf{Step 2: Resolve intended notional.} The target order size $Q_{\mathrm{intended}}$ is specified in USDT and converted to quote currency (BTC) using the indicative BTC/USDT rate at entry time. This normalization ensures that intended position sizes are comparable across BTC price regimes.

\textbf{Step 3: Track inventory and exit notional.} Let $R^{\mathrm{entry}}_{\mathrm{BTC/USDT}}$ denote the completed indicative conversion rate at entry. The entry quote notional and acquired base-asset quantity are
\begin{equation}
    Q^{\mathrm{entry}}_{\mathrm{BTC}}
      = Q_{\mathrm{intended}}/R^{\mathrm{entry}}_{\mathrm{BTC/USDT}},
    \qquad
    B = Q^{\mathrm{entry}}_{\mathrm{BTC}}/p^{\mathrm{entry}}_{\mathrm{vwap}}.
\end{equation}
The exit liquidates all $B$ units. Its impact is evaluated at the contemporaneous raw quote notional $Q^{\mathrm{exit}}_{\mathrm{BTC}}=B p^{\mathrm{exit}}$, rather than reusing the entry BTC notional. Internally the impact functions use USDT-normalized sizes $\widetilde Q^{\mathrm{entry}}=Q^{\mathrm{entry}}_{\mathrm{BTC}}R^{\mathrm{entry}}_{\mathrm{BTC/USDT}}$ and $\widetilde Q^{\mathrm{exit}}=Q^{\mathrm{exit}}_{\mathrm{BTC}}R^{T}_{\mathrm{BTC/USDT}}$, where $R^{T}$ is the completed announcement-time rate used to fit the manipulated exit model. This prevents exit slippage from being understated when the target price rises. We assume that deeper layers of the order book remain available beyond the volume observed in historical candles. Larger trades therefore remain feasible, but they incur larger slippage through the impact model.

\textbf{Step 4: Compute VWAP execution prices.} The fitted coefficient estimates the price displacement associated with a given net order flow. An order that walks through the book continuously fills at improving prices along the way. Integrating the square-root impact profile over the execution path yields the average fill cost:
\begin{equation}
    I_{\mathrm{vwap}}(Q)
      = \frac{1}{Q}\int_{0}^{Q}\!\beta\sqrt{v}\,dv
      = \tfrac{2}{3}\,\beta\sqrt{Q}.
    \label{eq:vwap_correction}
\end{equation}
The VWAP impact is two-thirds of the terminal impact because early fills occur at better prices than the final fill level. The VWAP execution prices are then:
\begin{align}
    p^{\mathrm{entry}}_{\mathrm{vwap}} &= p^{\mathrm{entry}}\;\bigl(1 + I_{\mathrm{vwap}}(\widetilde Q^{\mathrm{entry}}) / 10^4\bigr), \label{eq:vwap_entry}\\
    p^{\mathrm{exit}}_{\mathrm{vwap}}  &= p^{\mathrm{exit}}\;\bigl(1 - I_{\mathrm{vwap}}(\widetilde Q^{\mathrm{exit}}) / 10^4\bigr). \label{eq:vwap_exit}
\end{align}
Buying pushes the fill price above the initial quote; selling pushes it below. Let $R^{\mathrm{exit}}_{\mathrm{BTC/USDT}}$ be the completed exit-time conversion rate. The fully USDT-denominated transaction return is
\begin{equation}
    r_{\mathrm{USDT}}
      = \frac{B p^{\mathrm{exit}}_{\mathrm{vwap}} R^{\mathrm{exit}}_{\mathrm{BTC/USDT}}}
             {Q_{\mathrm{intended}}}
        - 1 - 0.0025,
    \label{eq:tx_return}
\end{equation}
where the 25 bps term is charged against entry capital and covers the round-trip exchange fee. This accounting includes the BTC/USDT move over the holding interval and never mixes a BTC-denominated asset return with an exit-time USDT conversion.

\textbf{Manipulated-asset exit.} The sell leg occurs during the manipulation window where buying and selling pressure are not balanced. Its impact model uses 5-second sell-dominated candles available strictly before the exit, never observations after the simulated transaction. A manipulation-window fit is used only when at least 20 sell-dominated candle samples are available; otherwise the framework logs the sample count and falls back to the pre-pump entry model.

\subsubsection{Regression fit quality}
Fig.~\ref{fig:impact_regression} shows an example of the fitted entry impact regression for one pump asset. The absolute-impact approach pools buy and sell 5-minute candles into a single regression ($I = \beta\sqrt{Q}$, no intercept), yielding $R^2 = 0.310$ on the buy side, $R^2 = 0.355$ on the sell side, and $\beta = 0.7224$ in the example shown.

Fig.~\ref{fig:exit_impact_regression} is a descriptive fit over the full capped 10-minute manipulation window for one asset. This longer diagnostic window is used only to visualize the sell-side relation; the execution backtest instead truncates each fit at that transaction's realized exit. The sell-only 5-second candle approach yields $R^2 = 0.274$ with $\beta = 1.3188$ ($n = 83$ candles) in the example shown. The higher $\beta$ reflects the increased cost of selling during a pump, consistent with the asymmetric liquidity conditions.

\begin{figure*}[!t]
    \centering
    \includegraphics[width=\textwidth]{images/impact_regression_example.png}
    \caption{Fitted entry price-impact regression for the SKY-BTC pair, using the 14 trading days of 5-minute candles ending 15 minutes before the pump announcement. Left panel: candles with net buying pressure (positive $Q_{\mathrm{net}}$); right panel: candles with net selling pressure (negative $Q_{\mathrm{net}}$). Each scatter point is one 5-minute candle; horizontal axis $|Q_{\mathrm{net}}|$ is absolute net volume in USDT and vertical axis $|\Delta p|$ is the absolute close-to-close price change in basis points. Solid line: fitted $|\Delta p| = \beta\sqrt{|Q_{\mathrm{net}}|}$ (no intercept, $\beta = 0.72$), estimated on both panels jointly. \textbf{Takeaway:} the single-parameter square-root fit captures the concave impact$-$size relation on both sides of the book, consistent with the square-root impact law documented across asset classes~\cite{toth2011anomalous,donier2015million}; this is the curve used to translate order size into VWAP slippage in the backtest.}
    \label{fig:impact_regression}
\end{figure*}

\begin{figure}[!t]
    \centering
    \includegraphics[width=\columnwidth]{images/exit_impact_regression_example.png}
    \caption{Descriptive exit price-impact regression for the DLT-BTC pair, using 5-second candles from the first 10 minutes after the pump announcement and restricted to sell-dominated candles ($n = 83$). This full-window illustration is not the execution fit: each backtested exit uses only candles strictly preceding its own exit timestamp. Horizontal axis: absolute sell-side net notional $|Q_{\mathrm{net}}|$ in USDT. Vertical axis: absolute close-to-close price change in basis points. Solid line: fitted $|\Delta p| = \beta\sqrt{|Q_{\mathrm{net}}|}$ ($\beta = 1.32$). \textbf{Takeaway:} the exit coefficient is nearly twice the entry coefficient ($1.32$ vs.\ $0.72$ on the earlier example), which reflects the asymmetric liquidity conditions during a pump---selling into the dying pump is materially more expensive than buying against resting depth pre-announcement.}
    \label{fig:exit_impact_regression}
\end{figure}

\subsubsection{Unmodeled execution frictions}
The fitted impact model captures average slippage against historical candle-level liquidity but intentionally omits several real-world frictions in order to keep the estimate interpretable. Specifically, we do not model (i) order-book depth exhaustion at the deepest levels (which would manifest as partial fills rather than a smooth $\beta\sqrt{Q}$ curve), (ii) queue priority and time-in-force dynamics on the exchange matching engine, (iii) exchange-side API latency between signal generation and order arrival, and (iv) the risk that other market participants detect the same pump signature and front-run the entry. Each of these frictions works in the same direction and would reduce realized PnL further, so the sensitivity numbers reported below should be read as an upper bound on achievable performance.

\subsubsection{PnL sensitivity to order size}
We sweep the intended order size from 100 to 10{,}000 USDT and re-run the impact-aware backtest for a $k=5$ portfolio using a 15-minute pre-announcement decision and a 1-minute post-announcement exit. Fig.~\ref{fig:pnl_vs_quantity} reports the no-reinvestment cumulative ROE over the full 80-event test sample as a function of intended order size under the candle-based impact model. Cumulative ROE declines monotonically from $0.840$ at 100 USDT to $-0.201$ at 10{,}000 USDT, remaining positive at 5{,}000 USDT ($0.135$) and crossing zero between 5{,}000 and 10{,}000 USDT. The strategy is therefore viable only at relatively small notionals under this execution model: beyond several thousand USDT per trade the square-root impact incurred on both legs consumes the edge. Larger positions would require more sophisticated execution algorithms (e.g., time-weighted average price (TWAP)/VWAP splitting), while institutional notionals lie further down the same adverse impact curve. Learning an optimal exit point instead of using the fixed 1-minute exit is another possible extension, but both directions require deeper order-book reconstruction than our candle-based fit supports.

\begin{figure}[!t]
    \centering
    \includegraphics[width=\columnwidth]{images/pnl_vs_quantity_with_impact.png}
    \caption{Cumulative return on invested capital (ROE) on the test sample as a function of intended per-trade order size, with the fitted candle-based square-root impact model enabled on both legs. Horizontal axis (log scale): intended order size in USDT, swept from $100$ to $10{,}000$. Vertical axis: cumulative ROE across the 80 test-set events for a $k=5$ portfolio with a 15-minute pre-announcement decision and full inventory liquidation at the 1-minute post-announcement exit. \textbf{Takeaway:} cumulative ROE falls monotonically from $0.840$ at $100$ USDT per trade to $-0.201$ at $10{,}000$ USDT and crosses zero between $5{,}000$ and $10{,}000$ USDT.}
    \label{fig:pnl_vs_quantity}
\end{figure}

\section{Limitations and Future Work}
\label{sec:limitations}

\textbf{Exchange generalizability.}
Our analysis is restricted to Binance BTC pairs. P\&D dynamics may differ on smaller exchanges with lower liquidity, different fee structures, or weaker surveillance. Validating the framework on additional exchanges is a natural extension, though it requires comparable data availability.

\textbf{Dataset size.}
With 473 labeled events (80 in the test set), statistical power is limited. The bootstrap confidence intervals (Section~\ref{sec:significance}) quantify this uncertainty, but larger labeled datasets would enable more precise performance estimation.

\textbf{Execution assumptions.}
Although we incorporate a fitted price-impact model estimated from pre-pump 5-minute candles and post-pump 5-second sell candles (Section~\ref{sec:price_impact}), the backtesting framework does not model all real-world frictions, including exchange API latency, order-book dynamics during the pump itself, or the risk of front-running by other participants who detect the same signal.

\textbf{Evolving manipulation tactics.}
Manipulators may adapt their strategies over time. Our temporal subperiod analysis (Section~\ref{sec:robustness}) provides evidence of cross-regime stability, but ongoing retraining and monitoring would be necessary for deployment.

\textbf{Alternative imbalance-handling approaches.}
While we compared class weighting and SMOTE, other techniques such as focal loss~\cite{lin2017focal}, cost-sensitive learning~\cite{elkan2001cost}, and deep anomaly detection~\cite{liu2008isolation} remain unexplored in this setting and could potentially improve performance.

\textbf{Feature audit.}
During the preparation of this revision we discovered a look-ahead in the z-score normalizers that computed the standardization moments on windows extending past the pump-time cut-off; we fixed the normalizer, tightened the feature cutoff from $T-1$ hour to $T-15$ minutes to match the portfolio entry time, re-anchored the hourly normalizer bars to the cutoff itself so every bar is a full hour of strictly-pre-cutoff trades, regenerated every feature parquet, re-tuned all hyperparameters, and recomputed every reported number in this paper from the leak-free rb-anchored features.


\section{Conclusion}

In this study, we successfully applied machine learning techniques to predict pump-and-dump targets, achieving strong empirical performance with our formulated approaches. Our analysis demonstrates that identifying P\&D targets in advance is feasible and presents opportunities for profitable trading strategies, albeit only at retail order sizes once the fitted price impact is charged on both legs (Section~\ref{sec:price_impact}). Specifically, the validation-selected CatBoost classifier with Top@k\%AUC early stopping achieved the highest test F1 and balanced accuracy among the evaluated models and supported an execution-aware portfolio strategy; tuned CatBoost achieved the highest PR-AUC. The selected model's aggregate test Top@k\%AUC exceeded that of the validation-selected tuned-CatBoost comparator by $0.0116$, but the pre-specified paired bootstrap test did not distinguish that difference from sampling noise at the 5\% level (Section~\ref{sec:significance}).

Our findings also explored alternative classification methods on imbalanced data, including ranking algorithms. However, these approaches demonstrated inferior performance when compared to boosting algorithms, such as the CatBoost classifier. Our research has significant implications for market surveillance and regulatory policy, highlighting the potential benefits of using machine learning techniques to detect and prevent pump-and-dump schemes.


\begin{thebibliography}{00}

\bibitem{optuna_2019} T.~Akiba, S.~Sano, T.~Yanase, T.~Ohta, and M.~Koyama, ``Optuna: A next-generation hyperparameter optimization framework,'' in \emph{Proc. 25th ACM SIGKDD Int. Conf. Knowl. Discovery Data Mining}, 2019.

\bibitem{albers_frag_2022} J.~Albers, M.~Cucuringu, S.~Howison, and A.~Y.~Shestopaloff, ``Fragmentation, price formation and cross-impact in Bitcoin markets,'' \emph{Appl. Math. Finance}, vol.~28, 2022.

\bibitem{random_forest_2001} L.~Breiman, ``Random forests,'' \emph{Mach. Learn.}, vol.~45, no.~1, 2001.

\bibitem{chadalapaka_crypto_2022} V.~Chadalapaka, K.~Chang, G.~Mahajan, and A.~Vasil, ``Crypto pump and dump detection via deep learning techniques,'' \emph{arXiv}, 2022.

\bibitem{cong_washtrading_2023} L.~W.~Cong, X.~Li, K.~Tang, and Y.~Yang, ``Crypto wash trading,'' \emph{Manage. Sci.}, vol.~69, no.~11, 2023.

\bibitem{econometrics11030022} D.~Fantazzini and Y.~Xiao, ``Detecting pump-and-dumps with crypto-assets: Dealing with imbalanced datasets and insiders' anticipated purchases,'' \emph{Econometrics}, vol.~11, no.~3, 2023.

\bibitem{grinsztajn_2022} L.~Grinsztajn, O.~Oyallon, and G.~Varoquaux, ``Why do tree-based models still outperform deep learning on typical tabular data?,'' in \emph{Advances in Neural Information Processing Systems}, 2022.

\bibitem{hu2023sequence} S.~Hu, Z.~Zhang, S.~Lu, B.~He, and Z.~Li, ``Sequence-based target coin prediction for cryptocurrency pump-and-dump,'' \emph{Proc. ACM Manag. Data}, vol.~1, no.~1, 2023.

\bibitem{kamps_moon_2018} J.~Kamps and K.~Bennett, ``To the moon: Defining and detecting cryptocurrency pump-and-dumps,'' \emph{Crime Sci.}, vol.~7, p.~18, 2018.

\bibitem{la_morgia_pump_2020} M.~La~Morgia, A.~Mei, F.~Sassi, and J.~Stefa, ``Pump and dumps in the Bitcoin era: Real-time detection of cryptocurrency market manipulations,'' in \emph{Proc. Int. Conf. Comput. Commun. Netw. (ICCCN)}, 2020, pp.~1--9.

\bibitem{la2023doge} M.~La~Morgia, A.~Mei, F.~Sassi, and J.~Stefa, ``The doge of wall street: Analysis and detection of pump-and-dump cryptocurrency manipulations,'' \emph{ACM Trans. Internet Technol.}, vol.~23, no.~1, pp.~1--28, 2023.

\bibitem{nghiem_2021} H.~Nghiem, G.~Muric, F.~Morstatter, and E.~Ferrara, ``Detecting cryptocurrency pump-and-dump frauds using market and social signals,'' \emph{Expert Syst. Appl.}, art. no. 115284, 2021.

\bibitem{ntakaris_midprice_2022} A.~Ntakaris, J.~Kanniainen, M.~Gabbouj, and A.~Iosifidis, ``Mid-price prediction based on machine learning methods with technical and quantitative indicators,'' \emph{PLOS ONE}, vol.~15, no.~6, 2020.

\bibitem{poh_ltr} D.~Poh, B.~Lim, S.~Zohren, and S.~Roberts, ``Building cross-sectional systematic strategies by learning to rank,'' \emph{J. Financial Data Sci.}, vol.~3, no.~2, pp.~70--86, 2021.

\bibitem{catboost} L.~Prokhorenkova, G.~Gusev, A.~Vorobev, A.~V.~Dorogush, and A.~Gulin, ``CatBoost: Unbiased boosting with categorical features,'' in \emph{Advances in Neural Information Processing Systems}, 2018.

\bibitem{pd_anatomy} J.~Xu and B.~Livshits, ``The anatomy of a cryptocurrency pump-and-dump scheme,'' in \emph{Proc. 28th USENIX Secur. Symp. (USENIX Security 19)}, 2019, pp.~1609--1625.

\bibitem{chawla2002smote} N.~V.~Chawla, K.~W.~Bowyer, L.~O.~Hall, and W.~P.~Kegelmeyer, ``SMOTE: Synthetic minority over-sampling technique,'' \emph{J. Artif. Intell. Res.}, vol.~16, pp.~321--357, 2002.

\bibitem{code_github}
M.~Mironov, D.~Bogutsky, and P.~Lukianchenko, ``Pump-and-dump target prediction (code repository), GitHub, 2025. [Online]. Available: https://github.com/B0R0koko/pump-and-dump-prediction. Accessed on: Dec.~24,~2025.


\bibitem{almgren2001optimal} R.~Almgren and N.~Chriss, ``Optimal execution of portfolio transactions,'' \emph{J. Risk}, vol.~3, no.~2, pp.~5--39, 2001.

\bibitem{almgren2005direct} R.~Almgren, C.~Thum, E.~Hauptmann, and H.~Li, ``Direct estimation of equity market impact,'' \emph{Risk}, vol.~18, no.~7, pp.~58--62, 2005.

\bibitem{blagus2013smote} R.~Blagus and L.~Lusa, ``SMOTE for high-dimensional class-imbalanced data,'' \emph{BMC Bioinform.}, vol.~14, p.~106, 2013.

\bibitem{bouchaud2009markets} J.-P.~Bouchaud, J.~D.~Farmer, and F.~Lillo, ``How markets slowly digest changes in supply and demand,'' in \emph{Handbook of Financial Markets: Dynamics and Evolution}, T.~Hens and K.~R.~Schenk-Hopp\'{e}, Eds.\hspace{0.25em} Elsevier, 2009, ch.~2.

\bibitem{donier2015million} J.~Donier and J.~Bonart, ``A million metaorder analysis of market impact on the Bitcoin,'' \emph{Market Microstructure and Liquidity}, vol.~1, no.~2, art.~1550008, 2015.

\bibitem{elkan2001cost} C.~Elkan, ``The foundations of cost-sensitive learning,'' in \emph{Proc. 17th Int. Joint Conf. Artif. Intell. (IJCAI)}, 2001, pp.~973--978.

\bibitem{kyle1985continuous} A.~S.~Kyle, ``Continuous auctions and insider trading,'' \emph{Econometrica}, vol.~53, no.~6, pp.~1315--1335, 1985.

\bibitem{lin2017focal} T.-Y.~Lin, P.~Goyal, R.~Girshick, K.~He, and P.~Doll\'{a}r, ``Focal loss for dense object detection,'' in \emph{Proc. IEEE Int. Conf. Comput. Vis. (ICCV)}, 2017, pp.~2980--2988.

\bibitem{liu2008isolation} F.~T.~Liu, K.~M.~Ting, and Z.-H.~Zhou, ``Isolation forest,'' in \emph{Proc. IEEE Int. Conf. Data Mining (ICDM)}, 2008, pp.~413--422.

\bibitem{toth2011anomalous} B.~T\'{o}th, Y.~Lemperiere, C.~Deremble, J.~de~Lataillade, J.~Kockelkoren, and J.-P.~Bouchaud, ``Anomalous price impact and the critical nature of liquidity in financial markets,'' \emph{Phys. Rev. X}, vol.~1, art.~021006, 2011.

\end{thebibliography}

\begin{IEEEbiography}[{\includegraphics[width=1in,height=1.25in,clip,keepaspectratio]{images/mironov.png}}]{Mikhail Mironov}
received the B.S.\ degree in economics from Universitat Pompeu Fabra, Barcelona, Spain, and the National Research University Higher School of Economics (HSE), St.\ Petersburg, Russia, in 2023 (double degree), and the M.S.\ degree with honors in financial economics from International College of Economics and Finance, ICEF HSE, Moscow, Russia, in 2025. He is currently pursuing the Ph.D.\ degree in finance at HSE, Moscow, Russia. His major field of study is finance and financial economics.

\end{IEEEbiography}

\begin{IEEEbiography}[{\includegraphics[width=1in,height=1.25in,clip,keepaspectratio]{images/bogutsky.png}}]{Denis Bogutsky} received the M.S. degree in mathematics from Lomonosov Moscow State University in 2020. He is currently a Junior Researcher with the Laboratory of AI in Mathematical Finance at the National Research University Higher School of Economics. His research focuses on quantitative finance, market microstructure, credit risk modeling, and anomaly detection in financial systems.
\end{IEEEbiography}

\newpage

%If you do not have or do not want to include a photo, you can use IEEEbiographynophoto as shown below:

\begin{IEEEbiography}
[{\includegraphics[width=1in,height=1.25in,clip,keepaspectratio]{images/lukyanchenko.png}}]
{PETR P. LUKIANCHENKO} received the bachelor’s and master’s degrees in information science from HSE University, Moscow, Russia, in
2009 and 2011, respectively, and the Ph.D. degree
in mathematics. He has been with HSE University,
since 2009.
He currently holds multiple positions, including
a Senior Lecturer with the Faculty of Computer
Science, Big Data and Information Retrieval
School; a Lecturer with the Centre for Training
Top AI Professionals; and the Head of the Laboratory of Artificial Intelligence in Mathematical Finance, AI and Digital Science Institute. At Trusted
AI Research Center, he conducts research projects with a particular focus
on image processing and time series data. He has co-authored research on
neural-network-based approaches for predicting anomalies in interest rates
under stochastic noise. He has presented his work on a multiagent simulator
for financial crisis modeling at AI Journey 2025 Conference. His research
interests include artificial intelligence applications in finance, including the
detection of structural breaks in financial time series and the development of
simulators for financial market crises.

\end{IEEEbiography}

\EOD

\end{document}
